\documentclass[11pt]{article}
\usepackage[T1]{fontenc}
\usepackage[utf8]{inputenc}
\usepackage{lmodern}
\usepackage[margin=0.86in]{geometry}
\usepackage{amsmath,amssymb,amsthm,mathtools}
\usepackage{booktabs,longtable,array,multicol}
\usepackage{microtype,xcolor}
\usepackage{enumitem}
\usepackage[hyphens]{url}
\usepackage{hyperref}
\hypersetup{colorlinks=true,linkcolor=blue!45!black,citecolor=blue!45!black,urlcolor=blue!45!black,
 pdftitle={History-dependent returns and a shallower 215-state ZFC machine},
 pdfauthor={Research contribution for ProveIt},pdfsubject={Exact Turing machines, consistency search, and source-calibrated compilation}}
\setlength{\parindent}{0pt}
\setlength{\parskip}{5pt}
\setlist{nosep,leftmargin=1.7em}
\newtheorem{theorem}{Theorem}[section]
\newtheorem{lemma}[theorem]{Lemma}
\newtheorem{proposition}[theorem]{Proposition}
\newtheorem{corollary}[theorem]{Corollary}
\theoremstyle{definition}
\newtheorem{definition}[theorem]{Definition}
\newcommand{\ZFC}{\mathsf{ZFC}}
\newcommand{\Con}{\operatorname{Con}}
\newcommand{\Halt}{\operatorname{Halt}}
\newcommand{\blank}{\mathsf{blank}}
\newcommand{\Nat}{\mathbb N}
\newcommand{\Kbase}{K_{215}}
\newcommand{\Kctx}{K_{215}^{\mathrm{hist}}}
\newcommand{\Ksh}{K_{215}^{\mathrm{shallow}}}
\newcommand{\code}[1]{\texttt{#1}}
\newcommand{\pin}{\texttt{75716498437bfc7291d5900cab2983a11fd54bac}}
\title{History-dependent returns and a shallower\newline 215-state ZFC machine\\[4pt]
\large Exact transition certificates, a three-command cache calculus,\newline and semantic refutation alarms}
\author{Research contribution for the ProveIt repository}
\date{21 October 2026}

\begin{document}
\maketitle
\begin{abstract}
We give two explicit deterministic, two-symbol, blank-tape machines with
215 working states, including initialization, and an additional terminal
state. The first changes five transitions of the pinned accepted K215
machine. Six short, context-dependent blocks prove its blank halting
equivalence by positive-duration stuttering. The useful observation is that
the preceding one or two transitions have already written the marker bits
needed to choose an existing return path. This releases two rows from the
register-engine obligation without adding a state.

The second machine combines those returns with a different decoder
allocation and a closed source reassociation. It has maximum source-address
depth 46 rather than 47, changes 1,042 primitive addresses, and has a
transition graph that is not a relabelling of the accepted table. Both
machines pass fresh quantified runtime derivations for all 6,176 guarded
primitive cases and independent complete blank-prefix checks. Using the
repository's inherited numerical-calculus calibration, ordinary reasoning
formalizable in ZFC gives
\(\neg\Halt(M,\blank)\leftrightarrow\Con(\ZFC)\) for each machine.

On the calculus side, we give a complete proof compiler for a three-command
transfer/cache design with 201 distinct source subtrees, and a sound
extension of the falsity alarm to every membership-atom code. These are
separate source-level results. No 214-state machine is claimed. The
remaining allocation obstruction and the precise scope of the negative
searches are recorded so that the next reduction can build on concrete
evidence.
\end{abstract}
\newpage
\tableofcontents
\newpage

\section{Results, conventions, and the pinned starting point}

The starting point is the repository's accepted
\path{Computability/BusyBeaver/Research/ZFC/candidate-k215} at commit
\begin{center}\small\pin.\end{center}
Its ordinary proof identifies blank halting with a finite ZFC refutation
\cite{k215}. The exact native table has SHA-256
\begin{center}\footnotesize\ttfamily
9683d6ed9628e6ce407c315b08f5da3d063da94e92706f275ca0f33e8272e4b2.
\end{center}
This input theorem and the inherited numerical calibration are credited
inputs, not new results of this packet. The new transition proofs, source
reassociation, exact tables, runtime runs, and source-calculus compiler
composition are supplied here.

\begin{table}[ht]
\centering\small
\begin{tabular}{@{}lrrrr@{}}
\toprule
Construction & TM states & Source subtrees & Leaves & Depth\\
\midrule
Pinned \(\Kbase\) & 215 & 202 & 4,751 & 47\\
History-return \(\Kctx\) & 215 & 202 & 4,751 & 47\\
Shallower decoder \(\Ksh\) & 215 & 202 & 4,751 & 46\\
Direct transfer/cache source & --- & 201 & 4,709 & 44\\
Membership-alarm source & --- & 206 & 4,746 & 47\\
\bottomrule
\end{tabular}
\caption{Accepted physical counts and separately measured source counts.
A dash means that no accepted physical machine of the target size is supplied
for that source. Source subtrees include distinct primitive leaves.}
\end{table}

The machine model is a two-way infinite tape indexed by \(\mathbb Z\), with
alphabet \(\{0,1\}\), initially all zero. A row \((q,b)\mapsto(w,d,z)\)
writes \(w\), moves one cell in direction \(d\in\{L,R\}\), and enters
\(z\). Native start is 0 and native halt is 254. All 215 working labels
have both read rows; halt has none. The nonconsecutive native labels are
inherited names, not extra states. Dense versions use working labels
\(0,\ldots,214\) and halt 215. Initialization is inside each counted table;
there is no externally supplied program tape or uncounted preprocessing.

The source consists of a full binary tree whose leaves are increments
\(\mathsf{CI}_j\), guarded decrements \(\mathsf{CD}_j\), structured breaks
\(\mathsf{CB}_k\), structured continues \(\mathsf{CC}_k\), and halt
\(\mathsf{CH}\). The nine logical registers have order
\((N,T,W,A,D,H,X,Y,L)\). A source-subtree count is a measure of shared
program syntax. It is not the binary state count: physical states can share
decoder and register-engine duties, and some decoder roles depend on a
known preceding tape bit.

The principal files are \path{machine/k215-shallow.native.tm} and
\path{machine/k215-shallow.tm}. The simpler bridge is
\path{machine/k215-context.native.tm}. Its small delta makes the new engine
argument particularly easy to audit. The source tables, complete compressed
runtime cases, exact startup traces, independent checkers, mathematical
notes, and inherited notices are included. A manifest pins every file.

\section{Using transition history as an already available marker}
\label{sec:returns}

\subsection{The five literal row changes}

Define \(\Kctx\) by changing exactly the following five rows of \(\Kbase\).
All labels in the proof below are native labels.
\begin{center}
\begin{tabular}{@{}rcll@{}}
\toprule
State & Read & Old write/move/target & New write/move/target\\
\midrule
13  & 1 & \(1,L,31\)   & \(1,L,81\)\\
27  & 1 & \(0,L,40\)   & \(0,L,302\)\\
83  & 1 & \(0,L,42\)   & \(0,L,57\)\\
156 & 1 & \(1,R,176\)  & \(1,R,52\)\\
198 & 0 & \(1,L,40\)   & \(1,L,54\)\\
\bottomrule
\end{tabular}
\end{center}
The new native table has SHA-256
\begin{center}\footnotesize\ttfamily
da20914ad9bda3f75de0381f2c05a7b5e67d26199b57b870a5ea3b72511ab627.
\end{center}
No label is added or removed. All 430 rows remain defined, and all working
states are reachable in the underlying directed transition graph.

The argument does not assert that these two tables act alike from every
arbitrary configuration. Their modified entries are equivalent on contexts
forced by actual incoming transitions. That is sufficient for runs from
the fixed start state, and it exposes sharing opportunities that a
context-free state minimizer misses.

\subsection{The exact incoming-context lemmas}

Write \(a_i\) for the tape symbol \(i\) cells from the head at entry.
Every entry to state 27 is through \(95/0=0R27\), so \(a_{-1}=0\).
Every entry to state 83 is through \(73/1=1R83\), so \(a_{-1}=1\).
Every entry to state 156 is through \(189/1=1L156\), so \(a_{1}=1\).
Every entry to state 198 is through \(66/0=1R198\), so \(a_{-1}=1\).
These statements hold in both tables. None of those entry states is the
start state.

The state-13 entry needs a two-transition suffix. Its possible incoming
edges are \(8/1=1R13\), \(38/0=0R13\), and \(42/0=0R13\). Enumerating all
two-edge incoming histories in each literal table establishes
\begin{equation}
 q=13\ \land\ a_0=1
 \quad\Longrightarrow\quad
 a_{-1}=1\ \lor\ (a_{-1}=0\land a_{-2}=1).
 \label{eq:history}
\end{equation}
The finite checker retains the latest write at each position while walking
the two edges backwards. It deliberately includes some read-incompatible
histories: this is an overapproximation, so checking the displayed property
on all retained histories is sound. An actual two-step suffix is always
represented. State 13 is neither the initial state nor a possible
one-transition successor of it, so there is no unhandled short-prefix
exception.

This use of history costs no machine state. The marker is already on the
tape because the incoming edge wrote it. The proof, rather than a new
runtime flag, carries the knowledge of that marker.

\subsection{Six complete local blocks}

The table lists the exact state paths. The entry bit and context columns
give every cell that the corresponding trace reads. No other tape cell is
inspected, so arbitrary contents outside the listed cells are allowed.

\begin{center}\small
\begin{tabular}{@{}llp{3.5cm}p{3.8cm}r@{}}
\toprule
Entry & Context & Old state path & New state path & Final head\\
\midrule
\(13/1\)& \(a_{-1}=0,a_{-2}=1\)&
\(13\to31\to21\)&\(13\to81\to111\to128\to21\)&0\\
\(13/1\)& \(a_{-1}=1\)&
\(13\to31\to47\to16\)&\(13\to81\to57\to16\)&\(-1\)\\
\(27/1\)& \(a_{-1}=0\)&
\(27\to40\to4\)&\(27\to302\to4\)&\(-2\)\\
\(83/1\)& \(a_{-1}=1\)&
\(83\to42\to16\)&\(83\to57\to16\)&\(-2\)\\
\(156/1\)& \(a_{1}=1\)&
\(156\to176\to150\)&\(156\to52\to150\)&2\\
\(198/0\)& \(a_{-1}=1\)&
\(198\to40\to54\)&\(198\to54\to54\)&\(-2\)\\
\bottomrule
\end{tabular}
\end{center}

In the first case both tapes are unchanged; the new trace takes four
transitions instead of two. In the second case both erase the entry cell
and leave the predecessor one unchanged. The third case erases the entry
one. The fourth erases both adjacent ones. The fifth changes no tape cell.
The sixth changes the entry zero to one. In every case the final state,
head, and entire tape are identical. Every strict interior transition of
every listed block is outside the set of five modified rows.

\begin{theorem}[Exact history-return equivalence]
From the fixed blank start,
\[
 \Halt(\Kctx,\blank)\quad\Longleftrightarrow\quad
 \Halt(\Kbase,\blank).
\]
The equivalence follows from a finite, positive-duration stuttering
simulation and is formalizable in ZFC.
\end{theorem}
\begin{proof}
Couple the two runs at identical configurations. If their next row is
unmodified, take the same single transition. Otherwise the incoming-context
lemma selects one of the six blocks. Both blocks are defined, do not halt
internally, and finish at an identical configuration. Their strict
interiors contain no modified row, so no nested replacement obligation is
being assumed. Iterate.

All block durations are positive and at most four. Thus infinitely many
coupled blocks cannot represent a finite run on either side, and a finite
halting prefix on either side yields one on the other. The entry-history
facts hold independently in both literal tables, including after earlier
replacements. This proves both directions without assuming consistency of
ZFC or eventual nonhalting.
\end{proof}

The independent checker \path{code/verify_context.py} reads the two frozen
tables directly. It checks the exact delta, all incoming guards, all local
traces, arbitrary exterior preservation, strict-interior separation, full
state count, canonical dense renaming, and the complete finite blank
prefix. It imports neither the allocator nor the symbolic runtime.

\subsection{What the engine change buys}

A fresh symbolic census of engine execution uses 142 transition rows,
compared with 144 for the accepted table. The removed obligations are
\(31/0\) and \(42/1\); no previously unused engine row is added. This is a
statement about engine-phase use, not a statement that these rows are
unreachable in the full machine. They still have decoder roles.

The count is established by rederiving every primitive case and separating
source-address reading from engine execution at each dispatcher visit.
A primitive may revisit the dispatcher internally, so simply treating all
rows after the first dispatch as engine rows gives an incorrect census.
The included census records use the corrected phase distinction.

Freeing an engine obligation enlarges the allowed decoder allocation. It
does not by itself save a working state. Conversely, keeping the old
destination of a row that the engine no longer uses would conceal this
opportunity from an allocator that freezes whole states.

\section{A different 215-state decoder}
\label{sec:shallow}

\subsection{One exact source rewrite}

Let the accepted source tree be \(S\). At its path
\begin{center}\ttfamily 11010010011101010110011100\end{center}
the subtree is \(((A,B),C)\). Replace this single occurrence by
\((A,(B,C))\). The factors contain 539, 380, and 123 primitive occurrences,
respectively. Each has zero escaping break and continuation demand and
contains no halt. Thus this is a reassociation of three closed sequential
blocks. It changes no primitive opcode, register operand, or control depth.

The emitted source is pinned by SHA-256
\begin{center}\footnotesize\ttfamily
c85089b913dcf446da5039843a43c15129efe88c8ee47a673b5b74600ae42ace.
\end{center}
Its source graph is checked twice: first by address arithmetic, including
the no-stray-suffix condition on each successor, and independently by
subtree intervals and ancestor unwinding. Both resolvers produce exactly
the accepted ordered effects and guarded successor ordinals, including
the initial ordinal and the halt leaf.

\begin{lemma}[Identity simulation of the source]
For every tuple of nine natural register values, the accepted source and
the reassociated source have identical register effects and branch targets
at every matching primitive ordinal. In particular they have the same
halting behavior from every matching source configuration.
\end{lemma}
\begin{proof}
Both source graphs are finite and literally equal. At an increment,
decrement, break, continue, or halt ordinal, the effect and each guarded
target agree. Induction on execution length proves the identity
simulation. No bound on a register value is used.
\end{proof}

There remain 202 distinct source subtrees and 4,751 primitive leaves.
Exactly 1,042 leaf addresses change, their total address length decreases
by 416, and maximum address depth decreases from 47 to 46. Maximum zero
run remains 6. These are literal source measurements, not estimates of
wall-clock speed or of the time before any possible refutation.

\subsection{The frozen physical table}

We allocate the rewritten source against the history-return engine, using
an engine-row filter and known-prefix checks. The resulting table
\(\Ksh\) has SHA-256
\begin{center}\footnotesize\ttfamily
7dafe445beb74f66b24d4e540136b5a47d0dfc6c1b877dcb8edce1526a608537.
\end{center}
It has 215 working states and 430 rows. The allocator changes 191 literal
rows relative to the pinned table. Part of this large row delta is decoder
relabelling, so the number alone would not establish a meaningful change.
The substantive differences are the history-dependent engine returns and
the changed source-address tree.

An independent finite graph check proves that the new table is not a
rooted, read-labelled, write/move-preserving state renaming of the original.
For example, the abstract read-label word \code{000000} gives different
write/move actions in the two rooted transition graphs. This is a graph
witness, not a claim that those reads form a blank-tape execution. A dense
canonical table and its complete native-to-dense bijection are supplied.

The allocator's prefix acceptance is only a synthesis result. The physical
table is accepted here because of the separate exact runtime and startup
checks described next. The proof does not infer binary correctness from
the source-node count, the MILP result, or a finite sample of register
values.

\section{Complete blank starts and quantified binary runtime}
\label{sec:runtime}

Both new tables reach the following complete calibrated frame after
95,740 transitions from a blank tape:
\[
\begin{gathered}
q=84,\qquad \text{head}=-8,\qquad \text{header}=57,\qquad \text{gap}=1,\\
(N,T,W,A,D,H,X,Y,L)=(4,0,0,0,0,0,0,0,0),\\
\text{spares}=[1]^{40}.
\end{gathered}
\]
The tape has exactly 95 one cells and is zero everywhere else. This is the
same complete frame as the accepted machine, whose prefix takes 95,714
transitions. There are 13 visits to the slower state-13 context in that
prefix, accounting exactly for the additional 26 transitions.

The independent frame reconstruction uses the physical order
\((W,L,T,D,A,Y,N,H,X)\). There is a one marker at \(-9\), a header one at
57, then unary tokens \(0\,1^{v+1}\) for the nine physical register values
and forty spare values. This independently gives 95 one cells, with last
one at 199. The checker reconstructs this set from the advertised
register frame and separately replays every recorded transition; equality
of the sets establishes the full zero exterior, not just a displayed tape
window.

The source seed has zero formula, stack, and cache registers. Its next
candidate value is 4. The inherited calibration expressly permits such a
finite cutoff: the proof compiler has cofinally many successful candidate
numbers whenever a refutation exists. The frame is therefore calibrated
without pretending that the enumerator starts at candidate zero.

\begin{table}[ht]
\centering\small
\begin{tabular}{@{}lrrr@{}}
\toprule
Verification quantity & Pinned table & \(\Kctx\) & \(\Ksh\)\\
\midrule
Working states & 215 & 215 & 215\\
Guarded primitive cases & 6,176 & 6,176 & 6,176\\
Symbolic inference steps & 1,698,790 & 1,699,668 & 1,697,385\\
Minimum constant duration & 6 & 6 & 6\\
Blank-prefix transitions & 95,714 & 95,740 & 95,740\\
\bottomrule
\end{tabular}
\caption{Counts of proof work and exact startup execution. The symbolic
inference totals are not TM running times.}
\end{table}

The fresh runtime derivations use the repository's archived affine-symbolic
checker, with the native table supplied explicitly. Each source increment,
break, continue, and halt contributes one guard case; each decrement has
zero and positive cases. The 4,751 occurrences therefore give 6,176 cases.
Each derivation quantifies over all natural primary register values
consistent with its guard, an arbitrary natural gap, and a nonempty
arbitrary finite spare list, including an arbitrary first spare and tail.

The proof rules consist of literal table transitions, induction over unary
scan lengths, and four separately derived list-block schemas. There is no
fixed maximum register value, proof length, or spare-list length in those
claims. The complete normalized frame and zero exterior are part of the
primitive contracts. Every reported affine duration has a positive
constant term and nonnegative variable coefficients; the minimum constant
is 6. Thus an unbounded source run corresponds to an unbounded binary run.

The exact table and source hashes, the checker and helper hashes, case
identities, and failures list are recorded in the runtime receipts. All
cases succeed. The independent shallow-decoder audit additionally checks
the complete unique guard coverage and binds the receipt to the included
runtime code. The compressed case lists and startup traces are included
without discarding difficult or exceptional entries.

These are ordinary mathematical derivations carried out by executable
checkers and independently reviewed arguments. They are not newly compiled
Lean or Rocq proofs of these exact tables. The packet does not claim a
proof-assistant kernel result, a PA endpoint, or an unconditional proof of
nonhalting.

\section{The resulting consistency equivalences}

\begin{theorem}[The supplied 215-state machines]
For each \(M\in\{\Kctx,\Ksh\}\), ordinary reasoning formalizable in ZFC
establishes
\begin{align*}
 \Halt(M,\blank)&\quad\longleftrightarrow\quad\neg\Con(\ZFC),\\
 \neg\Halt(M,\blank)&\quad\longleftrightarrow\quad\Con(\ZFC).
\end{align*}
Here consistency is the same arithmetized absence of a finite ZFC
refutation used in the pinned source calibration.
\end{theorem}
\begin{proof}
For \(\Kctx\), Section~\ref{sec:returns} proves blank halting equivalence
to the accepted exact table. Compose with that table's calibrated theorem
\cite{k215,calibration,dominated}. The fresh runtime and blank checks supply
an additional direct audit of the same implementation.

For \(\Ksh\), the independent blank replay reaches the calibrated source
root. The quantified primitive-runtime theorem relates every finite source
step to a finite positive-duration binary block. The equal source graph of
Section~\ref{sec:shallow} identifies that source with the calibrated
accepted source on all register values. Its command soundness invariant,
complete proof-block compiler, cofinal padding, and strictly draining
candidate reader yield halting exactly when the dominated numerical
calculus proves its falsity code. The inherited calibration identifies that
event with a finite ZFC refutation. The reverse simulation and positive
durations exclude an extra binary infinite behavior inside a source
primitive. The finite initial prefix then gives the blank-input statement.
\end{proof}

The consistency statement is not assumed in this proof of an equivalence.
If ZFC is consistent, the theorem has the expected unprovability
consequence; it does not license claiming that ZFC proves its own
consistency. All mathematical results needed beyond the new finite machine
and source arguments are identified in the cited repository calibration.

\begingroup
\newcommand{\ctP}{\mathsf P}
\newcommand{\ctI}{\mathsf I}
\newcommand{\ctQ}{\mathsf Q}
\newcommand{\ctD}{\mathcal D}
\newcommand{\ctS}{\mathsf S}
\newcommand{\ctMP}{\mathsf{MP}}
\newcommand{\ctAX}{\mathsf{AX}}
\newcommand{\ctH}{\mathsf H}
\newcommand{\ctPush}{\mathsf{PUSH}}
\newcommand{\ctDiscard}{\mathsf{DISCARD}}
\newcommand{\ctBlock}{\mathcal B}

\section{A 201-subtree source: transfer-pop and cache-return inference}
\label{sec:cache-calculus}

This construction combines the transfer-pop proof enumerator
of~\cite{transfer} with cache-return modus ponens from~\cite{cachemp},
rebased onto the accepted K215 source~\cite{k215}. The combination needs
a new proof compiler: cache-return inference no longer supplies the
dummy stack entry consumed when its result is transferred to the current
formula. Reserving that entry before compiling the premises resolves
the mismatch. The resulting literal programs have 201 distinct source
subtrees, one fewer than the accepted source. This section proves their
ordinary source-level consistency calibration. It does not identify a
214-state binary machine or infer a physical state count from the
source count; allocation, full-tape simulation, and blank initialization
are separate obligations.

\subsection{Numerical calculus and exact command interface}

Use Cantor pairing and the unchanged implication and quantifier codes:
\[
 \ctP(a,b)=\frac{(a+b)(a+b+1)}2+b,\qquad
 \ctI(a,b)=4\ctP(a,b),\qquad
 \ctQ_i a=4\ctP(i,a)+2.
\]
Truth is code0 and falsity is code1. In particular,
$\ctI(0,0)=0$. Write $[x;w]=\ctP(x,w)$ and use right-nested
pairing for longer stacks. This is numerical notation, not an
injective encoding of finite lists: $[0;0]=0$, and both projections
of zero are zero.

Let $\ctD$ be the accepted dominated numerical calculus: truth,
numerical MP, theorem-only generalization, the seven retained logical
rows $R_1,\ldots,R_7$, the complete gated theory bundle, and
\[
 L(A,B,C,U)=
 \ctI\bigl(\ctI(\ctI(C,U),A),
             \ctI(\ctI(A,C),\ctI(B,C))\bigr),\qquad U\leq A.
\]
These rows and the full bundle are retained literally. In particular,
none of the bundle's support or atomic-vacuity material is removed.
The inherited finite propositional certificate recovers unrestricted
$\rho_8$ for every natural code; the accepted calibration then gives
\begin{equation}
 \ctD\vdash1\quad\Longleftrightarrow\quad
 \exists p\;\operatorname{Proof}_{\mathrm{ZFC}}(p,\bot).
 \label{eq:ct-inherited}
\end{equation}
The present theorem uses this existing mathematical
calibration~\cite{k215}; it is not a new proof-assistant verification
of that dependency. Formula decoding and well-formedness tests are
not added to the numerical calculus.

The proof enumerator has a new operational language with unchanged
axiom strength. In particular, the operand of MP remains an
arbitrary natural increment; inference is not restricted to a fixed
consequent or a fixed parameter. Successful inference, cache movement,
and stack restoration are reorganized together. Treating either command
change in isolation would miss the obstruction solved by the reserved
zero, and would not justify transferring the old completeness theorem.

At command boundaries the work tuple is $(t,w,d,h)$: current formula,
paired stack, additive parameter, and cached formula. If
$\operatorname{unpair}(w)=(x,v)$, the direct transfer command is
\[
 \ctS(t,w,d,h)=
 \begin{cases}
  \mathsf{halt},&h=1,\\
  (h,v,d,x),&h\ne1.
 \end{cases}
\]
For $Y=d+x$, cache-return inference has the exact equation
\[
 \ctMP(x)(t,w,d,h)=
 \begin{cases}
  (0,w,0,Y),&\ctI(t,Y)=h,\\
  (0,w,Y,0),&\ctI(t,Y)\ne h.
 \end{cases}
\]
The successful branch clears the parameter and leaves the conclusion
in cache; it does not push a dummy zero. The failed branch is unchanged.

The axiom command $\ctAX(c;x,y)$ is unchanged. Set
$U=d$, $A=d+x$, $B=h+y$, and $C=c$. Its result is
\begin{align*}
 t'&=G(A,B,c),& d'&=0,&h'&=0,\\
 w'&=[R_1,\ldots,R_7,L(A,B,c,U),\ctQ_U h,t;w].
\end{align*}
Here $G$ is the original full bundle when $c\ne A,B$, and zero
otherwise. The quantified entry uses the incoming cache, before adding
$y$. Every logical row is an implication code and hence differs from1;
the quantified entry is nonzero of residue2 modulo4. The gate is also
never1: an enabled bundle is implication-coded, and a disabled gate
is zero. These facts hold for arbitrary raw operands.

\subsection{Literal construction and eight variants}

In the accepted tree, replace the successful-MP subtree at
\path{11010010111} by $(\mathrm{CD}_4,
(\mathrm{CI}_5,\mathrm{CC}_2))$. Delete the DROP selector layer
by replacing \path{11010} with its current descendant
\path{11010010}. Finally replace the old NIP body at \path{111}
by a transfer body. The direct body is the closed composition
\[
 \mathsf{clearT};\mathsf{copyH};\mathsf{oldNIP}.
\]
The first part clears current; the second is the old DROP cache-copy
tail, preserving cache and clearing scratch. Old NIP then tests that
cache and pops the stack. Both parenthesizations implement precisely
the direct $\ctS$ equation.

Alternatively use
$\mathsf{oldDROP};\mathsf{oldNIP};\mathsf{oldNIP}$, abbreviated
DNN. Its nonhalting map is the same, but it also halts if incoming
$t=1$. It is therefore an early-halting refinement of direct transfer.
Each body has two parenthesizations, and each program has two small
selector orders, giving eight literal variants:
\begin{center}
\begin{tabular}{lrrrrr}
\toprule
Transfer body&Variants&Distinct subtrees&Leaves&Depth&Zero run\\
\midrule
Direct&4&201&4709&44&6\\
DNN&4&201&4752&44&6\\
\bottomrule
\end{tabular}
\end{center}
For the ordinary order, wire fields are
\[
 f(\ctS)=[0],\qquad f(\ctMP(x))=[1,x],\qquad
 f(\ctAX(c;x,y))=[c+2,y,x].
\]
The \texttt{swap01} order exchanges only the first two selectors:
$f(\ctS)=[1]$ and $f(\ctMP(x))=[0,x]$. AX fields stay fixed.
The first two selector decrements prove these partitions for all
natural fields: zero selects the first body, one the second, and
$c+2$ enters AX with residual $c$.

The independent source audit reconstructs all eight files, checks both
CFG resolvers, and verifies all 4565 primitive effects and internal
successors of AX. Its unique exterior edge still reaches the common
branch continuation. Thus AX is not incorrectly assumed closed when
its required outer continuation is supplied by the selector.

\subsection{Soundness on all raw candidate words}

Call a work tuple safe if current, cache, and every successive decoded
stack head are $\ctD$-theorems. No theoremhood condition is imposed
on $d$. AX preserves safety by its inherited axiom and generalization
contracts. Failed MP inserts truth; successful MP derives $Y$ from
$t$ and $\ctI(t,Y)$, places it in cache, and inserts truth in current.
Transfer only moves theorem codes. Its cache-one guard, or either guard
of DNN, therefore witnesses a theorem equal to falsity. Consequently
every actual halt from a safe seed yields $\ctD\vdash1$. This argument
quantifies over all natural command fields, including unintended or
nondecodable inputs; it does not presuppose a valid compiled proof word.

\subsection{Structural blocks recovered without DROP}

Put $D_0=\ctAX(0;0,0)$ and define
\begin{align*}
 \ctH&=D_0;\ctS^{11},\\
 \ctPush&=D_0;\ctS^9;\ctMP(0),\\
 \ctDiscard&=\ctMP(1);\ctAX(1;0,0);\ctS^{10};\ctMP(0).
\end{align*}
From $(t,[x;w],d,h)$, the first ten shifts of $\ctH$ yield
$(\ctQ_dh,[x;w],0,t)$. Their tests inspect only zero and the nine
generated nonfalsity entries. The eleventh shift halts exactly when
$t=1$, otherwise returning $(t,w,0,x)$. Thus $\ctH$ preserves
current while popping one old stack entry and resetting the parameter.
It intentionally does not test the original cache before AX erases it.

After nine shifts in $\ctPush$, current is the dominated row,
cache is $\ctQ_dh$, and stack is $[t;w]$. The final MP necessarily
fails because an implication code has residue0 and this cache has
residue2. Hence PUSH never halts and returns $(0,[t;w],0,0)$ for
every natural input tuple.

DISCARD starts from $(t,w,0,0)$. Its first MP fails since
$\ctI(t,1)>0$. AX has $U=A=c=1$ and gate zero. Ten shifts give
$(\ctQ_1 0,w,0,0)$; the final MP fails because
$\ctI(\ctQ_1 0,0)>0$. Thus DISCARD never halts and returns
$(0,w,0,0)$. Finally, from current, cache, and parameter zero, define
$J(0)$ to be empty and $J(u)=\ctMP(u)$ when $u>0$. The latter is
a forced failure setting $d=u$. Omitting $\ctMP(0)$ is essential:
the zero alias would make it succeed.

For DNN, the extra top test does not invalidate these exact auxiliary
claims: every tested gate and generated entry is nonfalsity. The old
top remains untested in PUSH, and its sole test in H is the intended
last shift.

These blocks use the axiom emitter as a source of temporary storage.
Its fixed ten-entry output preserves the old current formula behind nine
safe entries; repeated transfers can recover that saved value or leave
it on the stack without a primitive DROP command. The intermediate
logical rows need not be useful premises in the proof being compiled.
Their exact stack positions, their nonfalsity tags, and the quantified
entry's residue are sufficient. This distinction matters for compression:
a command that is expensive as an isolated arithmetic routine can supply
several structural operations through already shared source subtrees.
The resulting proof words can be very long. No execution-time advantage
is required for a consistency machine, whose specification asks only
whether a finite halting run exists.

\subsection{The strong theorem-block compiler}

For every derivable code $X$, construct a finite block $\ctBlock_X$
which, from \emph{every} $(t,w,0,0)$, either halts or returns exactly
$(X,w,0,0)$. Arbitrary incoming $t,w$ need not be safe. This universal
completeness contract is distinct from the safe-run soundness invariant.
Truth uses DISCARD. The axiom cases are
\begin{align*}
 R_j:\quad&\ctDiscard;\ctAX(c;a,b);\ctS^{j+1};\ctH^{9-j}
                    &&(1\leq j\leq7),\\
 L(a,b,c,u):\quad&\ctDiscard;J(u);\ctAX(c;a-u,b);\ctS^9;\ctH
                    &&(u\leq a),\\
 G(a,b,c):\quad&\ctDiscard;\ctAX(c;a,b);\ctH^{10}
                    &&(c\ne a,b).
\end{align*}
The indicated shifts select the desired row; the H repetitions preserve
it and remove exactly the remaining generated entries, ending with
saved zero in cache. The underlying stack is restored. All auxiliary
MP calls are forced failures, so cache-return success does not affect
these calculations.

Generalization uses
\[
 \ctBlock_{\ctQ_iX}=
 \ctBlock_X;\ctPush;J(i);\ctS;\ctAX(i;0,0);\ctS^{10}.
\]
After PUSH, setter, and transfer, the state is $(0,w,i,X)$. That
transfer tests incoming cache zero, not the newly popped $X$. AX has
zero gate and generates $\ctQ_iX$ immediately above saved zero;
ten shifts return $(\ctQ_iX,w,0,0)$. This includes $i=0$ and
$X=1$, subject only to permitted earlier halting.

\paragraph{MP with a protected zero.}
For premises $A$ and $\ctI(A,Y)$ define
\begin{multline*}
 \ctBlock_Y=\ctDiscard;\ctPush;
 \ctBlock_{\ctI(A,Y)};\ctPush;
 \ctBlock_A;\ctPush;\\
 \ctS;\ctS;\ctMP(Y);\ctS.
\end{multline*}
If no earlier halt occurs, its exact state trace is
\begin{align*}
 (t,w,0,0)&\longmapsto(0,[0;w],0,0)\\
 &\longmapsto(0,[\ctI(A,Y),0;w],0,0)\\
 &\longmapsto(0,[A,\ctI(A,Y),0;w],0,0)\\
 &\longmapsto(0,[\ctI(A,Y),0;w],0,A)\\
 &\longmapsto(A,[0;w],0,\ctI(A,Y))\\
 &\longmapsto(0,[0;w],0,Y)\\
 &\longmapsto(Y,w,0,0).
\end{align*}
The second transfer may halt if $A=1$, and the final transfer halts
if $Y=1$. Otherwise the displayed equations hold. Premise blocks
preserve their entire incoming stacks, which protects the reserved zero
until it compensates for the final pop. Without this reservation, that
pop would destroy an arbitrary original stack head. Pairing identities
validate the trace even at $A=Y=w=0$; no list freeness is assumed.

Induction on a finite derivation proves the strong contract. A proof
DAG can be expanded to a finite tree before compiling the major premise
first. DNN transfers the contract unchanged: runs agree with the direct
version until DNN possibly halts earlier. If $\ctBlock_1$ returns,
append H to detect its current falsity. Thus $\ctBlock_1;\ctH$
succeeds from every $(t,w,0,0)$.

\subsection{Cofinal successful candidates and the source theorem}

Let $F(fs,r)$ right-fold the command fields above residual $r$, and set
\[
 \operatorname{tail}([])=0,\qquad
 \operatorname{tail}(k::L)=1+F(f(k),\operatorname{tail}(L)).
\]
A continuation first decrements its positive proof-list pointer;
a fresh candidate decodes without that outer successor. Put
$R=\operatorname{tail}(\ctBlock_1;\ctH)$. For each natural $z$,
\[
 N_z=\ctP\bigl(2,\ctP(z,\ctP(1,R))\bigr)
\]
starts with $\ctAX(0;1,z)$ in either selector order. This resets
$d=h=0$ from every retained tuple. Its current and stack may be
arbitrary, precisely as allowed by the strong contract. Since
$N_z\geq z$, successful candidates are cofinal beyond every fixed
finite startup cutoff.

Every nonhalting positive continuation strictly decreases: after its
predecessor operation, each field reader takes a Cantor right projection
no larger than that predecessor. The natural-counter macros terminate.
Hence each nonhalting candidate drains, and sequential enumeration
reaches a successful $N_z$ unless it has already halted.

\begin{theorem}[Source consistency calibration]
For each of the eight literal transfer/cache programs, started with
the unchanged enumerator scheduler and a safe initialized frame,
halting is equivalent to $\ctD\vdash1$. Consequently, using
\eqref{eq:ct-inherited}, its nonhalting is equivalent in ZFC to
$\operatorname{Con}(\mathrm{ZFC})$.
\end{theorem}

\begin{proof}
Safe-run soundness gives the forward implication. A native refutation
gives the strong block and cofinally many successful candidates, proving
the reverse implication. These finite syntactic transformations and
natural-number arguments are formalizable in ZFC. A binary realization
requires its own exact runtime and initializer proofs.
\end{proof}

\paragraph{Open parity-alarm experiment.}
A separate 205-subtree, 5025-leaf source replaces emitted equality
$E(a,b)=4\ctP(a,b)+3$ by $\ctI(E(a,b),1)$, interprets raw
residue3 atoms as disequality, and tests cache parity with the four-leaf
scanner $(\mathrm{CD}_5,((\mathrm{CD}_5,\mathrm{CC}_2),\mathrm{CH}))$.
Constant assignments in an irreflexive model make both odd atom classes
false. However, the false witness $1=\operatorname{Rel}(0,0)$ adds
syntactic variable0 to the equality representation. Semantic validity
alone does not establish support-sensitive vacuity or the exact
formula/binder coupling in the gated bundle. A revised encoder and
reverse calibration remain to be proved; this experiment is not covered
by the theorem above.
\endgroup

\section{A semantic alarm: every membership atom detects inconsistency}
\label{sec:membership-alarm}

The halting test need not recognize one distinguished contradiction code.
It may recognize any decidable family of codes that contains that
contradiction and contains no theorem under the consistency hypothesis.
This observation separates the semantic choice of a halting predicate from
the cost of implementing it.  The following application changes the alarm
from equality with~$1$ to the congruence $h\equiv1\pmod4$.
It is an alternative source-level construction; the physical allocation
reported below does not produce a new 215-state machine.

\subsection{A general replacement principle}

Let $\mathcal D$ be the calibrated numerical calculus, with falsity code
$f$, and let $E$ be its refutation enumerator.  At command boundaries,
assume the usual invariant: the current formula, the cache, and every
decoded stack head are $\mathcal D$-theorems.  The original alarm halts
when the cache is~$f$.  Its other commands preserve this invariant, and
its proof compiler and scheduler establish refutation completeness.

\begin{proposition}[Semantic alarm replacement]
\label{prop:semantic-alarm}
Let $A\subseteq\mathbb N$ be decidable, with $f\in A$.  Suppose ZFC proves
\[
 \operatorname{Con}(\mathrm{ZFC})\ \Longrightarrow\
 \forall a\in A\;\neg(\mathcal D\vdash a).
\]
Replace the original alarm by membership in $A$, preserving its complete
nonhalting command effect.  From any calibrated seed, the modified
enumerator halts exactly when ZFC is syntactically inconsistent.
\end{proposition}

\begin{proof}
Under $\operatorname{Con}(\mathrm{ZFC})$, the invariant excludes an alarm
value at every command boundary.  Each nonhalting modified command has
the old invariant-preserving effect, so the new enumerator cannot halt.

For the converse, compare the two enumerators from the same raw boundary
state.  Until the new one halts, their boundary stores coincide.  When
the old alarm would halt, its cache is $f\in A$, so the new alarm also
halts, unless it has already halted earlier.  Induction on a finite old
halting execution proves that every such execution is detected by the
new enumerator.  Applying the old refutation-completeness theorem gives
the result.  This comparison does not assume that the starting store is
safe: its only premise is equality of the nonhalting command effects.
\end{proof}

A sufficient semantic condition is that, assuming consistency, each
$a\in A$ has a counterassignment in a model in which numerical
derivations are sound.  The model and assignment may depend on $a$.
The application below uses one nonempty model and one constant assignment
to refute every member of the alarm family simultaneously.

\subsection{All raw codes of residue one}

Write
\[
 P(i,j)=\frac{(i+j)(i+j+1)}2+j,
 \qquad R(i,j)=4P(i,j)+1.
\]
The total numerical decoder used in the inherited calibration satisfies
$\operatorname{dec}(R(i,j))=\operatorname{Rel}(i,j)$.
Its membership and soundness lemmas are \texttt{dec\_wel},
\texttt{sat\_wel}, and \texttt{c254\_sound} in
\texttt{Code254.lean}.\cite{code254}
Cantor unpairing gives, for every natural $h$,
\[
 h\equiv1\pmod4
 \quad\Longleftrightarrow\quad
 \exists i,j\in\mathbb N\;h=R(i,j).
\]
Indeed, unpairing $(h-1)/4$ supplies the required indices.  Thus every
accepted raw number is a membership atom.  No restriction to numbers
emitted by a proof compiler, no freshness test, and no separate
well-formedness test is present.  The truth alias~$0$ is excluded.

\begin{lemma}
In a nonempty irreflexive native model satisfying the calibrated bundles,
no membership-atom code is derivable in $\mathcal D$.
\end{lemma}

\begin{proof}
Let the model have domain $V$ and relation $\in_M$, with
$\neg(v\in_M v)$ for every $v\in V$.  Choose $v\in V$ and put
$e(k)=v$ for every natural variable index~$k$.
If $\mathcal D\vdash R(i,j)$, code soundness makes its decoded formula
true under every assignment, hence under~$e$.  The literal atom clause
then gives $v\in_M v$, a contradiction.  The argument applies also when
$i\ne j$: distinct variable names may receive the same value.
The dominated calculus is contained in the calibrated unrestricted
calculus, so its soundness follows directly; the stronger theorem-set
equality established for dominated $\rho_8$ is not needed in this
direction.\cite{dominated}
\end{proof}

The model-existence step is conditional and takes place within ZFC.
From the syntactic statement $\operatorname{Con}(\mathrm{ZFC})$, the
first-order completeness theorem supplies a set-coded model of ZFC.
The inherited native/ZFC interpretation then supplies a nonempty native
model satisfying the exact gated bundles.  Its membership relation is
irreflexive, as required by the calibrated semantics.\cite{k215,code254}
No transitive model and no externally well-founded model is assumed.
Foundation is used through what the model satisfies, not through an
external well-foundedness claim about its relation.
Consequently ZFC proves
\[
 \operatorname{Con}(\mathrm{ZFC})\Longrightarrow
 \forall h\in\mathbb N\,
   \bigl(\mathcal D\vdash h\Longrightarrow h\not\equiv1\pmod4\bigr).
\]
This is precisely the soundness premise of
Proposition~\ref{prop:semantic-alarm}, with $f=1$.

The actual-run argument covers unintended command encodings as well as
compiled proof words: the old invariant says that the cache is a theorem
at every alarm entry.  Conversely, $1\equiv1\pmod4$, so every old
refutation is still detected.  The strong compiler's ``halt or return
exactly'' contracts, its cofinal padding above arbitrary finite cutoffs,
and the decreasing continuation rank therefore transfer unchanged for
completeness.\cite{k215}
Expanding the alarm further to all odd numbers would be unsound for
this encoding: the equality code $4P(i,i)+3$ is derivable and has odd
parity.  The distinction between the two atomic tags is essential.

\subsection{Exact scanner and measured cost}

The residue test can destructively scan the cache with a four-phase
cycle.  Each phase tests for zero and otherwise decrements it.  The
zero exits of phases $0,1,2,3$ are respectively
\[
  \text{return},\qquad\text{halt},\qquad
  \text{return},\qquad\text{return}.
\]
Four positive decrements return to phase~$0$.
Induction on the quotient in $h=4q+r$ proves the exact behavior for
every natural cache value, and every normal exit has cache zero.
Appending the unchanged stack-pop tail gives precisely the old NIP
nonhalting output, including its scratch-register conditions.

The supplied search exhausts a declared finite family of seven-leaf
structured scanners: four cache decrements, one halt, one backward
continue, and one forward break, with closed control flow.
It checks their resolved phase graphs rather than testing finitely many
cache values.  Of 28,452 specifications, three have the required graph.
The best resulting complete source has 206 distinct subtrees,
4,746 primitive occurrences, maximum depth~47, and maximum zero run~6.
The exact source is \texttt{sites-atomhalt-best.json}; its search record
and reconstruction program accompany it.

The attempted 215-state allocation was infeasible in the declared
protected-row allocation model.  This is neither a lower bound for the
formalism nor a claim of a 206-state Turing machine: 206 counts source
subtrees.  No accepted 214- or 215-state binary implementation of this
alarm is supplied.  A future implementation must still prove its
all-natural primitive runtime and reach a calibrated blank-tape seed.
The semantic theorem remains independent of that unresolved encoding
cost, and no global priority claim is made for the replacement principle.


\section{What the searches establish, and what remains open}
\label{sec:search}

The source and physical searches were run against the pinned 215 table,
not against a stale larger candidate. Their useful output is both the
explicit alternatives above and a more specific account of why several
attractive source reductions do not immediately lower the binary count.

\subsection{Source size and physical sharing are different objectives}

The successful-MP implementation can be shortened in several ways.
Returning the result in the cache or omitting the dummy stack push each
gives a 202-subtree source at this stage. Combining cache return with the
three-command transfer system gives 201 subtrees. Nevertheless, ordinary
row-protected allocation at 215 states failed for the tested direct,
reassociated, and opcode-swapped versions. It also failed with the two new
engine freedoms from Section~\ref{sec:returns}.

One concrete reason is that an old successful-MP subprogram participates
in a context-sensitive marker fold. Removing its useful semantic work can
also remove its physical sharing opportunity. For the 201-subtree design,
the global allocation still needs 133 newly assigned decoder roles plus
32 fixed physical decoder states. Matching this total to the old total
does not guarantee compatibility with the surviving engine rows and
helper destinations.

A diagnostic relaxation identifies particular obstructions: allowing the
allocator to overwrite any one of the protected rows \(26/0\), \(26/1\),
\(37/0\), or \(302/0\) permits a 215-state prefix allocation of the direct 201-subtree
source. These relaxed outputs are \emph{not accepted machines}; they would
alter engine execution without a proof. Moving the engine's \(37/0\)
return to another existing provider did not by itself resolve the entire
allocation. This narrows a future attack to a concrete selector interface
instead of suggesting that one missing source node is automatically one
missing binary state.

\subsection{Finite search boundaries}

The recorded searches include 287 global source-reassociation candidates,
of which 216 passed the exact graph condition. Thirty-one distinct
202-subtree ties were saved; none in that family had 201 subtrees. A
cache-return-specific search considered 729 local reassociations with
structured-control-depth repair and obtained 137 valid graph candidates,
again without a 201-subtree improvement of that source. These are bounded
families of syntax transformations, not lower bounds over all source
programs.

On the physical side, 56,658 existing-helper context-fold specifications
were checked on the original table and 58,788 on the history-return table,
with no additional accepted fold. Another 26,896 ordered direct-context
contractions had no survivor. A local retargeting search considered
274,631 specifications and supplied the exact history-dependent blocks
used here. Its failure to find a state deletion does not imply that every
possible multirow redesign is impossible.

A final search synthesized both rows of an existing helper while
contracting a decoder pair. It emitted 336 structural 214-state trials;
none passed the exact symbolic block/primitive derivation filter. These
rejected trials are not candidate solutions. Separately, freeing all
admissible singleton nonpure decoder roles simultaneously gave a valid
215-state control allocation, but still no allocation for the direct
201-subtree source. The recorded negative results remain relative to their
specified synthesis and proof-rule families.

The membership-alarm source search enumerated 28,452 structured scanner
specifications with four cache decrements and one each of halt, continue,
and break. It accepted three exact four-phase scanner graphs. Its best
complete source had 206 distinct subtrees. The acceptance condition was an
exact finite control graph establishing the all-natural residue behavior,
not testing the first several hundred cache values.

There is also an exact zero/positive abstract-state analysis of the
accepted source. It identifies the always-zero decrement at the beginning
of the successful-MP dummy push and two consequent unreachable leaves.
Deleting that empty transfer alone produces 203 source subtrees: duplicated
surrounding structure offsets the local deletion. This is another example
of why raw leaf savings can be misleading.

\subsection{The next concrete targets}

The strongest uncompleted target is the three-command cache-return source.
It has a complete source-level consistency compiler and fewer distinct
subprograms, but its engine/decoder overlap still needs a successful exact
allocation. A useful next experiment is a multirow redesign of the
state-26 selector, paired with allocation of its old decoder role; a local
engine return that is correct only after entry through state 11 must be
checked under that specific entry contract and then subjected to the full
runtime checker.

The other target is a cheaper code family for a semantic alarm. The
membership class demonstrates that the alarm need not be a distinguished
literal falsity number. The challenge is to make the safe class cheap to
recognize without increasing the cost of all formula constructors or
silently breaking the source's numerical binder and support conditions.
The separate negative-equality parity note marks precisely that remaining
calibration obligation.

The two accepted machines in this packet are implementation alternatives
at 215 states. The 201-subtree calculus and semantic alarm are substantive
formalism directions. Neither is described as an achieved 214-state
breakthrough, and no exhaustive optimality or global novelty claim is
made. Existing cache-return and transfer-pop ideas are credited to the
integrated repository reports \cite{cachemp,transfer}. The newest incoming
binary ZIP submissions were not available as readable source through the
text retrieval interface, so priority comparisons are limited to the
integrated material actually inspected.

\clearpage
\section{Reproduction and incorporation}

Run the following from the extracted packet directory with Python 3.10 or
later. The finite checks and the symbolic runtime use the standard library.
The allocator alone needs NumPy and SciPy and is not needed to verify the
frozen machines.
\begin{verbatim}
python3 code/verify_context.py
python3 code/audit_shallow.py
python3 code/audit_runtime_receipts.py
python3 research/verify_cache_transfer.py
python3 research/build_cache_transfer.py
\end{verbatim}
The last command also reruns explicitly labelled finite command regressions.
The separate all-natural command proof is in this report and
\path{research/cache_transfer_proof.md}.

To rederive the unbounded primitive contracts rather than only checking
their archived receipts:
\begin{verbatim}
python3 code/runtime-inputs/validate_native.py \
  machine/k215-context.native.tm machine/sites215.json \
  reproduction/context-runtime

python3 code/runtime-inputs/validate_native.py \
  machine/k215-shallow.native.tm machine/sites-shallow.json \
  reproduction/shallow-runtime
\end{verbatim}
The exact checker code is vendored with its inherited notices. These runs
write complete startup traces and per-guard case records, so their output
is larger than a short success log. A failure to derive a case would be a
failed certificate, not automatically a counterexample to the machine's
semantics.

Compile the paper with \code{pdflatex report.tex} twice, or use the included
\path{Makefile}. The dense full transition table in the appendix is
generated from the frozen native shallow-decoder table by
\path{code/audit_runtime_receipts.py}. Thus the PDF also contains a complete
machine specification; the machine-readable tables remain the preferred
input for independent work.

The packet is intended as an additive research contribution. Inherited
tables, utilities, checker code, mathematical calibration, notices, and
attributions retain their existing identity and terms. The root README
states the exact new claims, source-only results, proof status, and file
layout. The archive contains no claim that the new exact tables have been
compiled through a proof-assistant kernel.

\clearpage
\begingroup\raggedright
\begin{thebibliography}{9}
\bibitem{k215}
ProveIt contributors, \emph{Ordinary proof of the combined215 machine},
\path{Computability/BusyBeaver/Research/ZFC/candidate-k215/proof.md},
commit \pin.
\href{https://github.com/VladimirReshetnikov/ProveIt/blob/75716498437bfc7291d5900cab2983a11fd54bac/Computability/BusyBeaver/Research/ZFC/candidate-k215/proof.md}{Pinned source}.

\bibitem{calibration}
ProveIt contributors, \emph{All-natural calibration of the permuted command
source}, candidate-k215 provenance, \path{permuted-commands/sections/source_calibration.tex},
same commit. A copy is included as \path{upstream/source_calibration.tex}.

\bibitem{dominated}
ProveIt contributors, \emph{Dominated numerical-calculus proof},
\path{Computability/BusyBeaver/Research/ZFC/candidate-k218/alternatives/dominated/proof.md},
same commit.
\href{https://github.com/VladimirReshetnikov/ProveIt/blob/75716498437bfc7291d5900cab2983a11fd54bac/Computability/BusyBeaver/Research/ZFC/candidate-k218/alternatives/dominated/proof.md}{Pinned source}.

\bibitem{code254}
ProveIt contributors, \emph{Code254.lean},
\path{Computability/BusyBeaver/Lean/BusyBeaver/ZFC/Code254.lean},
same commit; in particular the literal membership constructor, its
decoding/satisfaction lemmas, and code soundness.
\href{https://github.com/VladimirReshetnikov/ProveIt/blob/75716498437bfc7291d5900cab2983a11fd54bac/Computability/BusyBeaver/Lean/BusyBeaver/ZFC/Code254.lean}{Pinned source}.

\bibitem{cachemp}
ProveIt contributors, \emph{Cache-return modus-ponens proof}, integrated
report \path{54-cachemp-proofs-cache_mp_proof.md},
\path{SetTheory/Cardinals/docs/reports/set-theory/zfc-closure-kernels/},
same commit. The cache-return interface predates the present composition.

\bibitem{transfer}
ProveIt contributors, \emph{Complete proof compiler for the three-command
NIP-transfer enumerator}, integrated report
\path{166-transfer-proofs-transfer_compiler.md}, in the same report
directory and commit; see also its composite-transfer review.
\end{thebibliography}
\endgroup

\clearpage
\appendix
\section{Complete dense transition table for the shallower 215-state machine}

Start in state 0 on an all-zero two-way tape. The working states are
\(0,\ldots,214\); the only terminal state is 215. An entry \(wDz\)
means write \(w\), move in direction \(D\), and enter \(z\). Each line
below gives both read cases for three working states. The final empty right
block is only table layout, not a missing transition.

\begingroup\small
\begin{longtable}{@{}rll@{\hspace{1.4em}}rll@{\hspace{1.4em}}rll@{}}
\toprule
State & Read 0 & Read 1 & State & Read 0 & Read 1 & State & Read 0 & Read 1 \\
\midrule
\endfirsthead
\toprule
State & Read 0 & Read 1 & State & Read 0 & Read 1 & State & Read 0 & Read 1 \\
\midrule
\endhead
\midrule\multicolumn{9}{r}{\emph{Continued on the next page}}\\
\endfoot
\bottomrule
\endlastfoot
0 & \texttt{1R1} & \texttt{1R0} & 72 & \texttt{0R89} & \texttt{1R90} & 144 & \texttt{0R159} & \texttt{1R160} \\
1 & \texttt{0R2} & \texttt{0R0} & 73 & \texttt{1R73} & \texttt{1L13} & 145 & \texttt{0L161} & \texttt{1R162} \\
2 & \texttt{1R3} & \texttt{1L4} & 74 & \texttt{0R51} & \texttt{1R74} & 146 & \texttt{0R46} & \texttt{1L137} \\
3 & \texttt{1L4} & \texttt{1L5} & 75 & \texttt{0L73} & \texttt{0R91} & 147 & \texttt{0R98} & \texttt{1R163} \\
4 & \texttt{0L6} & \texttt{1L4} & 76 & \texttt{0R92} & \texttt{1R76} & 148 & \texttt{0R145} & \texttt{1R164} \\
5 & \texttt{0L7} & \texttt{0L8} & 77 & \texttt{0R46} & \texttt{1R93} & 149 & \texttt{0R111} & \texttt{1R52} \\
6 & \texttt{0L9} & \texttt{1L4} & 78 & \texttt{0R94} & \texttt{1R68} & 150 & \texttt{0R80} & \texttt{1R72} \\
7 & \texttt{0L9} & \texttt{0L7} & 79 & \texttt{0R95} & \texttt{1R96} & 151 & \texttt{0L164} & \texttt{1R105} \\
8 & \texttt{1L5} & \texttt{1L8} & 80 & \texttt{0R97} & \texttt{1R98} & 152 & \texttt{0L37} & \texttt{0R50} \\
9 & \texttt{1R10} & \texttt{0L9} & 81 & \texttt{0R44} & \texttt{1R99} & 153 & \texttt{0R89} & \texttt{1R121} \\
10 & \texttt{0R10} & \texttt{0R11} & 82 & \texttt{0L100} & \texttt{1L101} & 154 & \texttt{0R155} & \texttt{0L12} \\
11 & \texttt{0L12} & \texttt{0R11} & 83 & \texttt{0R102} & \texttt{1R103} & 155 & \texttt{0L53} & \texttt{1R165} \\
12 & \texttt{1L13} & \texttt{0L13} & 84 & \texttt{0R104} & \texttt{1R105} & 156 & \texttt{0R166} & \texttt{0L12} \\
13 & \texttt{0L13} & \texttt{1L14} & 85 & \texttt{0R106} & \texttt{1R107} & 157 & \texttt{0R105} & \texttt{1R78} \\
14 & \texttt{0L15} & \texttt{1L14} & 86 & \texttt{0R108} & \texttt{1R109} & 158 & \texttt{0R98} & \texttt{1L103} \\
15 & \texttt{0L16} & \texttt{1L14} & 87 & \texttt{0R75} & \texttt{0R92} & 159 & \texttt{0R78} & \texttt{1R83} \\
16 & \texttt{0L17} & \texttt{1L14} & 88 & \texttt{0R110} & \texttt{1R111} & 160 & \texttt{0R167} & \texttt{1R168} \\
17 & \texttt{0L18} & \texttt{1L14} & 89 & \texttt{0R57} & \texttt{0R92} & 161 & \texttt{0R126} & \texttt{1R169} \\
18 & \texttt{0L19} & \texttt{1L14} & 90 & \texttt{0R41} & \texttt{0L12} & 162 & \texttt{0R78} & \texttt{1R104} \\
19 & \texttt{0L20} & \texttt{1L14} & 91 & \texttt{0R112} & \texttt{1R91} & 163 & \texttt{0R80} & \texttt{1R125} \\
20 & \texttt{0L21} & \texttt{1L14} & 92 & \texttt{0R38} & \texttt{1R92} & 164 & \texttt{0R170} & \texttt{1R133} \\
21 & \texttt{0L22} & \texttt{1L14} & 93 & \texttt{0R113} & \texttt{0L7} & 165 & \texttt{0R53} & \texttt{1L40} \\
22 & \texttt{0R22} & \texttt{1R23} & 94 & \texttt{0R72} & \texttt{1R114} & 166 & \texttt{0R123} & \texttt{1R123} \\
23 & \texttt{0R24} & \texttt{1R25} & 95 & \texttt{0R115} & \texttt{1R44} & 167 & \texttt{0R84} & \texttt{1R158} \\
24 & \texttt{0R26} & \texttt{1R27} & 96 & \texttt{0R116} & \texttt{1R117} & 168 & \texttt{0R171} & \texttt{1R158} \\
25 & \texttt{0R28} & \texttt{1R29} & 97 & \texttt{0R33} & \texttt{1R118} & 169 & \texttt{0R151} & \texttt{1R172} \\
26 & \texttt{0R30} & \texttt{0L7} & 98 & \texttt{0R119} & \texttt{1R120} & 170 & \texttt{0R158} & \texttt{1R173} \\
27 & \texttt{0R31} & \texttt{1R32} & 99 & \texttt{0R121} & \texttt{0L12} & 171 & \texttt{0R104} & \texttt{1R78} \\
28 & \texttt{0R33} & \texttt{1R34} & 100 & \texttt{0R104} & \texttt{1R122} & 172 & \texttt{0R102} & \texttt{1R174} \\
29 & \texttt{0R35} & \texttt{1R36} & 101 & \texttt{0R123} & \texttt{1R124} & 173 & \texttt{0R175} & \texttt{1R176} \\
30 & \texttt{0L37} & \texttt{0R38} & 102 & \texttt{0R125} & \texttt{1R126} & 174 & \texttt{0R177} & \texttt{1R104} \\
31 & \texttt{0R39} & \texttt{1R40} & 103 & \texttt{0R127} & \texttt{1R128} & 175 & \texttt{0R98} & \texttt{1R178} \\
32 & \texttt{0R39} & \texttt{1R41} & 104 & \texttt{0R94} & \texttt{1R129} & 176 & \texttt{0R179} & \texttt{1R180} \\
33 & \texttt{0R42} & \texttt{1R43} & 105 & \texttt{0R94} & \texttt{1R130} & 177 & \texttt{0R82} & \texttt{1R181} \\
34 & \texttt{0R44} & \texttt{1R45} & 106 & \texttt{0R131} & \texttt{1R132} & 178 & \texttt{0R167} & \texttt{1R182} \\
35 & \texttt{0R46} & \texttt{1R47} & 107 & \texttt{0R133} & \texttt{0L7} & 179 & \texttt{0R183} & \texttt{1R181} \\
36 & \texttt{0R41} & \texttt{1R48} & 108 & \texttt{0R57} & \texttt{0R74} & 180 & \texttt{0R184} & \texttt{1R185} \\
37 & \texttt{0L29} & \texttt{1R49} & 109 & \texttt{0R87} & \texttt{1R56} & 181 & \texttt{0L186} & \texttt{1R187} \\
38 & \texttt{0R50} & \texttt{1R38} & 110 & \texttt{0R44} & \texttt{0L11} & 182 & \texttt{0R123} & \texttt{1R168} \\
39 & \texttt{0L37} & \texttt{0R51} & 111 & \texttt{0R56} & \texttt{1R90} & 183 & \texttt{0R188} & \texttt{1R189} \\
40 & \texttt{0R30} & \texttt{1L52} & 112 & \texttt{1R0} & \texttt{1R91} & 184 & \texttt{0L190} & \texttt{1L83} \\
41 & \texttt{0L53} & \texttt{1L54} & 113 & \texttt{0R37} & \texttt{1R134} & 185 & \texttt{0R181} & \texttt{1R146} \\
42 & \texttt{0R30} & \texttt{1R55} & 114 & \texttt{0R135} & \texttt{1R136} & 186 & \texttt{0R184} & \texttt{1R191} \\
43 & \texttt{0R56} & \texttt{1R40} & 115 & \texttt{0R108} & \texttt{1R117} & 187 & \texttt{0R82} & \texttt{1R98} \\
44 & \texttt{0R57} & \texttt{0R58} & 116 & \texttt{0L137} & \texttt{0L13} & 188 & \texttt{0R192} & \texttt{1R161} \\
45 & \texttt{0R59} & \texttt{0L12} & 117 & \texttt{0R110} & \texttt{0L7} & 189 & \texttt{0R190} & \texttt{1R193} \\
46 & \texttt{0R57} & \texttt{0R60} & 118 & \texttt{0R138} & \texttt{1R105} & 190 & \texttt{0R175} & \texttt{1R82} \\
47 & \texttt{0R61} & \texttt{0L7} & 119 & \texttt{0R139} & \texttt{1R140} & 191 & \texttt{0R184} & \texttt{1R194} \\
48 & \texttt{0R62} & \texttt{1R63} & 120 & \texttt{0R141} & \texttt{1R142} & 192 & \texttt{0R123} & \texttt{1R184} \\
49 & \texttt{0R64} & \texttt{1R65} & 121 & \texttt{0L42} & \texttt{1L86} & 193 & \texttt{0R127} & \texttt{1R195} \\
50 & \texttt{0R12} & \texttt{1R50} & 122 & \texttt{0R143} & \texttt{1R127} & 194 & \texttt{0R196} & \texttt{1R197} \\
51 & \texttt{0R66} & \texttt{1R51} & 123 & \texttt{0R87} & \texttt{1R87} & 195 & \texttt{0R198} & \texttt{1R186} \\
52 & \texttt{0R67} & \texttt{1L68} & 124 & \texttt{0R144} & \texttt{1R145} & 196 & \texttt{0R187} & \texttt{1R181} \\
53 & \texttt{0R69} & \texttt{1R54} & 125 & \texttt{0R105} & \texttt{1R104} & 197 & \texttt{0R199} & \texttt{1R196} \\
54 & \texttt{0L28} & \texttt{1R70} & 126 & \texttt{0R98} & \texttt{1R123} & 198 & \texttt{0R200} & \texttt{1R201} \\
55 & \texttt{0R69} & \texttt{1R71} & 127 & \texttt{0R126} & \texttt{1R98} & 199 & \texttt{0R184} & \texttt{1R202} \\
56 & \texttt{0L72} & \texttt{1L55} & 128 & \texttt{0R64} & \texttt{1R87} & 200 & \texttt{0R101} & \texttt{1R203} \\
57 & \texttt{0L73} & \texttt{0R0} & 129 & \texttt{0R146} & \texttt{1R34} & 201 & \texttt{0R143} & \texttt{1R204} \\
58 & \texttt{0R74} & \texttt{1R58} & 130 & \texttt{0R52} & \texttt{1R138} & 202 & \texttt{0R205} & \texttt{1R205} \\
59 & \texttt{0R75} & \texttt{0R60} & 131 & \texttt{0R63} & \texttt{1R147} & 203 & \texttt{0R206} & \texttt{1R177} \\
60 & \texttt{0R76} & \texttt{1R60} & 132 & \texttt{0R148} & \texttt{1R149} & 204 & \texttt{0R207} & \texttt{1R208} \\
61 & \texttt{0R77} & \texttt{1R78} & 133 & \texttt{0R79} & \texttt{1R116} & 205 & \texttt{0R162} & \texttt{1L209} \\
62 & \texttt{0R79} & \texttt{1R215} & 134 & \texttt{0R150} & \texttt{1R151} & 206 & \texttt{0R184} & \texttt{1R210} \\
63 & \texttt{0L80} & \texttt{1R81} & 135 & \texttt{0R152} & \texttt{1R70} & 207 & \texttt{0R100} & \texttt{1R211} \\
64 & \texttt{0L82} & \texttt{1L83} & 136 & \texttt{0R141} & \texttt{1R140} & 208 & \texttt{0R212} & \texttt{1R197} \\
65 & \texttt{0R84} & \texttt{1R85} & 137 & \texttt{0R86} & \texttt{1R153} & 209 & \texttt{0R145} & \texttt{1R168} \\
66 & \texttt{0R60} & \texttt{1R66} & 138 & \texttt{0R44} & \texttt{1R154} & 210 & \texttt{0R64} & \texttt{1R213} \\
67 & \texttt{0R57} & \texttt{0R76} & 139 & \texttt{0R152} & \texttt{1L54} & 211 & \texttt{0R171} & \texttt{1R205} \\
68 & \texttt{0R86} & \texttt{1R81} & 140 & \texttt{0R56} & \texttt{1R139} & 212 & \texttt{0R125} & \texttt{1R214} \\
69 & \texttt{0R75} & \texttt{0R66} & 141 & \texttt{0R153} & \texttt{1R155} & 213 & \texttt{0R102} & \texttt{1R64} \\
70 & \texttt{0R87} & \texttt{0L12} & 142 & \texttt{0R56} & \texttt{1R156} & 214 & \texttt{0R125} & \texttt{1R160} \\
71 & \texttt{0R88} & \texttt{0L70} & 143 & \texttt{0R157} & \texttt{1R158} &  & &  \\
\end{longtable}
\endgroup

\end{document}
